% ID: SEL-D0032-OB-E03
% ESTADO: cotejada. Segunda pasada visual y borrador compilado revisados.
% Fuente: Selectividad/matematicas_aplicadas/titular_junio_EXAMEN 5-ANDALUCÍA.pdf
% SHA-256: 5f24b200fb6cf34f6d72016ec3dc35787b4074de9b7752677e5b0d959ed0dfcd
% Página: 2; opción B; ejercicio original 3.

En una ciudad hay dos fábricas de pasta, F1 y F2, que producen dos tipos de productos, A y B, que venden a un distribuidor en paquetes de 1 kg. En un mes, la fábrica F1 produce 20000 kg de pasta, de los que 12000 son del tipo A y la fábrica F2 produce 25000 kg de pasta de los que 15000 kg son del tipo A. Se escoge al azar un paquete del distribuidor.
\begin{enumerate}[label=\alph*)]
\item \textbf{(1.5 puntos)} \textquestiondown{}Cuál es la probabilidad de que sea del tipo B?
\item \textbf{(1 punto)} Si el paquete elegido resulta ser del tipo A, \textquestiondown{}qué es más probable, que proceda de la fábrica F1 o que proceda de la F2?
\end{enumerate}
